% =================================================
% Часть 1. Подготовка к решению
% =================================================
\worksheettitle
  {Часть 1. Готовимся к решению}
  {Четыре ключевых перехода}
\studentnameblock



\section{Различаем три похожие записи}

\begin{fillbox}[Что означает число \(2\)?]
  \begin{center}
    \renewcommand{\arraystretch}{1.6}
    \begin{tabularx}{0.96\textwidth}{@{}
      >{\centering\arraybackslash}p{0.19\textwidth}
      >{\raggedright\arraybackslash}p{0.38\textwidth}X@{}}
      \toprule
      \rowcolor{TableHeader}
      Запись & Как прочитать & Роль числа \(2\) \\
      \midrule
      \(\lg^2 A\) & \answerblank[48mm] & \answerblank[44mm] \\
      \addlinespace[0.3em]
      \(\lg(A^2)\) & \answerblank[48mm] & \answerblank[44mm] \\
      \addlinespace[0.3em]
      \(\log_2 A\) & \answerblank[48mm] & \answerblank[44mm] \\
      \bottomrule
    \end{tabularx}
  \end{center}
\end{fillbox}

\notationtrianglefigure

\begin{checkpointbox}[Быстрая проверка]
  Вычислите:
  \[
    \lg^2 1000=\answerblank[22mm],
    \qquad
    \lg(1000^2)=\answerblank[22mm].
  \]
  Объясните словами, почему ответы различаются.
  \writinglines{2}
\end{checkpointbox}

\section{Преобразуем логарифм, не теряя ОДЗ}

\begin{fillbox}[Начинаем с области определения]
  Для выражения \(\lg((x-1)^2)\) аргумент должен быть положительным:
  \[
    (x-1)^2>0
    \quad\Longleftrightarrow\quad
    x\ne\answerblank[14mm].
  \]
  Подставьте \(x=0\). Какое выражение определено?
  \[
    \square\;\lg((x-1)^2)
    \qquad
    \square\;2\lg(x-1).
  \]
\end{fillbox}

\begin{fillbox}[Выводим формулу, не сужая область определения]
  При \(z\ne0\):
  \[
    z^2=|z|^2,
    \qquad
    \lg(z^2)=2\lg|z|.
  \]
  Поэтому
  \[
    \lg((x-1)^2)=2\lg\answerblank[26mm],
    \qquad x\ne1.
  \]
  Аналогично при \(z\ne0\):
  \[
    \lg(z^4)=\answerblank[12mm]\lg\answerblank[18mm].
  \]
\end{fillbox}

\begin{warningbox}[Что именно делает модуль]
  Модуль не появляется из-за правила «под квадратом всегда модуль».
  Он позволяет заменить \(z\) положительным числом \(|z|\), не изменяя
  чётную степень и не теряя отрицательные значения \(z\).
\end{warningbox}

\newpage
\section{Читаем модуль как расстояние}

\begin{methodintro}[Модуль как расстояние]
  Неравенство
  \[
    |u|\le3
  \]
  означает: расстояние от числа \(u\) до нуля не больше \(3\).
  Поэтому
  \[
    u\in\left[\answerblank[10mm],\answerblank[10mm]\right].
  \]
\end{methodintro}

\begin{center}
\begin{tikzpicture}[x=1.25cm,y=1cm]
  \draw[-{Stealth[length=2.2mm]}] (-3.7,0)--(3.9,0);
  \draw[very thick] (-3,0)--(3,0);
  \fill (-3,0) circle (2.4pt);
  \fill (3,0) circle (2.4pt);
  \foreach \x in {-3,0,3}{\draw (\x,0.10)--(\x,-0.10) node[below=4pt]{\(\x\)};}
\end{tikzpicture}
\end{center}

\begin{checkpointbox}[Изменяем только одно условие]
  Теперь решите
  \[
    0<|u|\le3.
  \]
  Условие \(|u|\le3\) оставляет отрезок \([-3;3]\), а условие
  \(|u|>0\) исключает единственное число \(u=\answerblank[12mm]\).
  Поэтому
  \[
    u\in
    \left[\answerblank[10mm],\answerblank[10mm]\right)
    \cup
    \left(\answerblank[10mm],\answerblank[10mm]\right].
  \]
\end{checkpointbox}

\begin{center}
\begin{tikzpicture}[x=1.25cm,y=1cm]
  \draw[-{Stealth[length=2.2mm]}] (-3.7,0)--(3.9,0);
  \draw[very thick] (-3,0)--(-0.11,0);
  \draw[very thick] (0.11,0)--(3,0);
  \fill (-3,0) circle (2.4pt);
  \draw[fill=white,very thick] (0,0) circle (2.7pt);
  \fill (3,0) circle (2.4pt);
  \foreach \x in {-3,0,3}{\draw (\x,0.10)--(\x,-0.10) node[below=4pt]{\(\x\)};}
\end{tikzpicture}
\end{center}

\begin{reflectionbox}[Сравните]
  Чем отличаются множества решений \(|u|\le3\) и \(0<|u|\le3\)?
  Сформулируйте ответ одним предложением.
  \writinglines{2}
\end{reflectionbox}

\section{Получаем два полезных правила}

\begin{fillbox}[Два условия выполняются одновременно]
  При \(a>0\):
  \[
    0<|u|\le a
    \quad\Longleftrightarrow\quad
    \begin{cases}
      |u|>0,\\
      |u|\le a.
    \end{cases}
  \]
  Поэтому
  \[
    0<|u|\le a
    \quad\Longleftrightarrow\quad
    -a\le u<0
    \quad\text{или}\quad
    0<u\le a.
  \]
  Второй полезный шаблон:
  \[
    |u|\ge a
    \quad\Longleftrightarrow\quad
    u\le-a
    \quad\text{или}\quad
    u\ge a.
  \]
\end{fillbox}

\begin{checkpointbox}[Применяем правило на простом примере]
  Решите:
  \[
    0<|x-2|\le3.
  \]
  Сначала положите \(u=x-2\), найдите промежутки для \(u\), затем вернитесь
  к \(x\).
  \workarea{20mm}
  \longanswerline
\end{checkpointbox}

\newpage
\section{Возвращаемся от квадрата к переменной}

\begin{fillbox}[Две симметричные части ответа]
  Решите неравенство
  \[
    1\le x^2<4.
  \]
  Сначала перепишите его через \(|x|\):
  \[
    \answerblank[12mm]\le|x|<\answerblank[12mm].
  \]
  Затем запишите промежутки для \(x\):
  \[
    x\in\answerblank[95mm].
  \]
\end{fillbox}

\begin{checkpointbox}[Соединяем два перехода]
  Решите неравенство
  \[
    0<|x^2-4|\le1.
  \]
  \begin{enumerate}
    \item Положите \(u=x^2-4\) и раскройте условие для \(u\).
    \item Получите два условия на \(x^2\).
    \item Перейдите к промежуткам для \(x\).
  \end{enumerate}
  \workarea{20mm}
  \longanswerline
\end{checkpointbox}

\begin{questionsbox}[Итоговая проверка перед частью 2]
  Отметьте только те утверждения, которые можете объяснить словами.
  \begin{itemize}[label=\choicebox,leftmargin=2.4em]
    \item Я различаю \(\lg^2 A\), \(\lg(A^2)\) и \(\log_2 A\).
    \item Я объясняю, почему \(\lg(z^2)=2\lg|z|\) при \(z\ne0\).
    \item Я вижу на числовой прямой решение \(0<|u|\le a\).
    \item Я перехожу от условий на \(x^2\) к обеим частям ответа для \(x\).
  \end{itemize}
\end{questionsbox}
